\paperheader
  {Michelson's Interference Experiment}
  {H. A. Lorentz}
  {Translated from "Versuch einer Theorie der elektrischen und optischen Erscheinungen in bewegten Körpern," Leidem, 1895, \S\S\,89-92.}


\section*{1}
As Maxwell first remarked and as follows from a very simple calculation, the time required by a ray of light to travel frome a point A to a point B and back to A must vary when the two points together undergo a displacement without carrying the ether with them. The differnce is, certainly, a magnitude of second order; but it is sufficiently great to be detected by a sensitive interference methiod.

The experiment was carried out by Michelson in 1881.\footnote{Michelson, American Journal of Science, 22, 1881, p. 120.} His apparatus, a kind of interferometer, had two hrizontal arms, P and Q, of equal length and at right angles one to the other. Of the two mutually interfering rays of light the one passed along the arm P and back, the other along the arm Q and back. The whole instrument, including the source of light and the arrangement for taking observations, could be recolved about a vertical axis; and those two positions come especially under consideration in which the arm P or the arm Q lay as nearly as possible in the direction of the Earth's motion. On the basis of Fresnel's theory it was anticipated that when the apparatus was revolved from one of these \textit{principal positions} into the other there would be a displacement of the interference fringes.

But of such a displacement - for the sake if brevity we will call it the Maxwell displacement - conditioned by the change in the times of propagation, no trace was discoverd, and accordingly Michelson thought himself justified in concluding that while the Earth is moving, the ether does not remain at rest. The correctness of this inference was soon brought into question, for by an oversight Michelson had taken the change in the phase differnce, which was to be expected in accordance with the theory, at twice its proper value. If we make the necessary correction, we arrive at displacements no greater than might be masked by errors of observation.

Subsequently Michelson \footnote{Michelson and Morley, American Journal of Science, 34, 1887, p. 333; Phil. Mag., 24, 1887, p. 449.} took up the investigation anew in collaboration with Morley, enhancing the delicacy of the experiment by causing each pencil to be reflected to and from between a number of mirrors, thereby obtaining the same advantage as if the arms of the earlier apparatus had been considerably lengthened. The mirrors wre mounted on a massive stone disc, floating on mercury, and therefore easily revolved. Each pencil now had to travel a total distance of 22 meters, and on Fresnel'stheory the displacement to be expected in passing from the one principal position to the other would be 0.4 of the distance between the interference fringes. Nevertheless the rotation produced displacements not exceeding 0.02 of this distance, and these might well be ascribed to errors of observation.

Now, does this result entitle us to assume that the ether takes part in the motion of the Earth, and therefore that the theory of aberration given by Stokes is incorrect one? The difficulties which this theory encounters in explaining aberration seem too great for me to share this opinion, and I would rather try to remove the contradiction between Frensel's theory and Michelson's result. An hypothesis which I brought forward some time ago, \footnote{Lorentz, Zittingsverslagen der Akad. v. Wet. te Amsterdam, 1892-93, p.74.} and which, as I subsequently learned, has also occurred to Fitzgerald, \footnote{As Fitzgerald kindly tells me, he has for a long time dealt with his hypothesis in his lectures. The only published reference which I can find to the hypothesis is by Lodge, "Aberration Problems," Phil. Trans. R.S., 184A, 1893.} enables us to do this. The next paragraph will set out this hypothesis.


\section*{2}
To simplify matters we will assume that we are working with apparatus as employed in the first experiment, and that in the one principal position the arm P lies exactly in the direction of the Earth's motion. Let $v$ be the velocity of this motion, $L$ the length of either arm, and hence $2L$ the path traversed by the rays of light. Accordin to the theory, \footnote{Cf. Lorentz. Arch. Néerl., 2, 1887, pp.168-176.}
